\section{灵敏度分析}

为检验模型的稳健性，对关键参数与假设进行扰动分析。

\subsection{问题一模型设定的稳健性}

在混合效应模型中分别加入孕周与 BMI 的交互项、孕周二次项，比较三种设定的系数与显著性。结果显示三种设定下孕周系数均为正、BMI 系数均为负且均显著，交互项与二次项不显著，说明``Y 浓度随孕周线性上升、随 BMI 线性下降''的结论对模型设定不敏感。

\subsection{问题二/三检测误差的灵敏度}

检测误差 $\sigma_e$ 从 0.005 增大到 0.020，最佳时点相应后移 2.0 周至 8.0 周（表~\ref{tab:q2_err}）。这表明最佳时点对检测精度敏感：检测误差增大时，为维持 90\% 的可靠达标率，必须显著推迟检测时点。因此，临床上提升检测精度可有效提前可靠检测时间，从而降低中期/晚期发现的风险。

\subsection{问题三分组方式的稳健性}

分别用仅 BMI 分组与多维聚类分组，二者给出的``高 BMI 组最佳时点更晚''的结论方向一致；多维聚类在组内同质性上更优（组内标准差之和 37.12 对比 40.44）。结论对分组方式不敏感。

\subsection{问题四判定阈值的灵敏度}

对基线 $|Z|>3$ 的阈值分别取 2.5 与 3.5 考察灵敏度--特异度权衡：阈值降低可提高灵敏度但降低特异度，阈值升高则相反；随机森林模型在 AUC 指标上始终稳定优于阈值法，说明多维综合判定对阈值选择更鲁棒。

综上，本文各模型的定性结论在参数扰动下保持稳健，主要定量结果（达标时间--BMI 关系、最佳时点、判定 AUC）均稳定。
